\section{General affine seminorm spectra}\label{sec:spectrum-foundations}

The affine seminorm spectrum in this section is defined for every commutative ring,
without choosing a coefficient field or a norm on the ring. Its points are real-valued
multiplicative ultrametric seminorms. For a prescribed seminorm on a coefficient ring,
the relative spectrum consists of the points restricting to that seminorm.
For an algebra over a trivially valued field this recovers the affine analytification
used in \cite[\S2.1]{FHX}. This construction does not impose the boundedness condition
that accompanies the spectrum of a Banach algebra.

\begin{definition}[Affine seminorm spectrum]\label{spectrum:affine-seminorm}
\lean{EnvelopeProperty.Berkovich.Spectrum, EnvelopeProperty.Berkovich.Spectrum.ofValuation}
\leanok
For a commutative ring $A$, let $\mathcal S(A)$ be the set of functions
$x\colon A\to\R_{\ge0}$ satisfying $|0|_x=0$, $|1|_x=1$,
$|ab|_x=|a|_x|b|_x$, and $|a+b|_x\le\max\{|a|_x,|b|_x\}$.
Every real-valued valuation on $A$ gives such a point by its underlying seminorm.
\end{definition}

\begin{lemma}[Pointwise topology]\label{spectrum:topology}
\lean{EnvelopeProperty.Berkovich.Spectrum.isEmbedding_toFun,
EnvelopeProperty.Berkovich.Spectrum.continuous_abs,
EnvelopeProperty.Berkovich.Spectrum.continuous_iff}
\leanok
\uses{spectrum:affine-seminorm}
Equip $\mathcal S(A)$ with the coarsest topology making every evaluation
$x\mapsto|a|_x$ continuous. It is Hausdorff. A map $g$ from a topological
space into $\mathcal S(A)$ is continuous if and only if every
$z\mapsto|a|_{g(z)}$ is continuous.
\end{lemma}
\begin{proof}
Embed $\mathcal S(A)$ into the product $\R_{\ge0}^A$ by listing all its values.
Two seminorms with the same list of values are equal, so this map is injective.
The induced topology has precisely the asserted characterization of continuity,
because a map into a product is continuous exactly when all its coordinates are.
A subspace of a Hausdorff product is Hausdorff.
\end{proof}

\begin{lemma}[Logarithmic evaluations]\label{spectrum:logarithmic-evaluation}
\lean{EnvelopeProperty.Berkovich.Spectrum.logAbs,
EnvelopeProperty.Berkovich.Spectrum.continuous_logAbs,
EnvelopeProperty.Berkovich.Spectrum.logAbs_mul,
EnvelopeProperty.Berkovich.Spectrum.logAbs_pow,
EnvelopeProperty.Berkovich.Spectrum.logAbs_add_le,
EnvelopeProperty.Berkovich.Spectrum.logAbs_eq_bot_iff,
EnvelopeProperty.Berkovich.Spectrum.logAbs_ne_top}
\leanok
\uses{spectrum:topology}
The function $\log|a|(x)=\log|a|_x$ is continuous with values in
$[-\infty,+\infty)$, equals $-\infty$ exactly when $|a|_x=0$, and satisfies
$\log|ab|=\log|a|+\log|b|$, $\log|a^n|=n\log|a|$, and
$\log|a+b|\le\max\{\log|a|,\log|b|\}$.
\end{lemma}
\begin{proof}
The extended logarithm is a continuous order-preserving map
$[0,+\infty]\to[-\infty,+\infty]$. Composing it with evaluation gives continuity.
Each evaluation is finite, so its logarithm is never $+\infty$; the logarithm is
$-\infty$ precisely at zero. Applying the product, power, and maximum identities
for this logarithm to the seminorm axioms gives the asserted formulas.
\end{proof}

\begin{lemma}[Pullback of seminorms]\label{spectrum:pullback}
\lean{EnvelopeProperty.Berkovich.Spectrum.comap,
EnvelopeProperty.Berkovich.Spectrum.comap_val_apply,
EnvelopeProperty.Berkovich.Spectrum.continuous_comap,
EnvelopeProperty.Berkovich.Spectrum.comap_id,
EnvelopeProperty.Berkovich.Spectrum.comap_comp}
\leanok
\uses{spectrum:topology}
For a unital ring homomorphism $\varphi\colon A\to B$, restriction defines
a continuous map $\varphi^*\colon\mathcal S(B)\to\mathcal S(A)$ with
$|a|_{\varphi^*y}=|\varphi(a)|_y$. Restriction preserves identity maps and
reverses composition.
\end{lemma}
\begin{proof}
Composing a seminorm with a unital ring homomorphism preserves each seminorm axiom.
Its evaluation at $a$ is evaluation at $\varphi(a)$ before restriction, hence is continuous.
The continuity criterion therefore gives continuity of restriction.
For the identity homomorphism the evaluation at every $a$ is unchanged.
For $A\xrightarrow{\varphi}B\xrightarrow{\psi}C$, both restrictions evaluate
$a$ as $|\psi(\varphi(a))|$, proving equality of the two maps.
\end{proof}

\begin{definition}[Relative spectrum]\label{spectrum:relative}
\lean{EnvelopeProperty.Berkovich.relativeLocus,
EnvelopeProperty.Berkovich.RelativeSpectrum,
EnvelopeProperty.Berkovich.isClosed_relativeLocus,
EnvelopeProperty.Berkovich.RelativeSpectrum.self,
EnvelopeProperty.Berkovich.RelativeSpectrum.isClosedEmbedding_subtypeVal,
EnvelopeProperty.Berkovich.RelativeSpectrum.continuous_comap}
\leanok
\uses{spectrum:pullback}
Let $R$ be a commutative ring, $v\in\mathcal S(R)$, and $A$ a commutative
$R$-algebra. The relative spectrum $\mathcal S_v(A)$ is the subspace
of $\mathcal S(A)$ consisting of the seminorms with
$|r|_x=|r|_v$ for every $r\in R$. It is closed in $\mathcal S(A)$.
An $R$-algebra homomorphism $A\to B$ induces a continuous restriction
$\mathcal S_v(B)\to\mathcal S_v(A)$.
The prescribed seminorm itself is a point of $\mathcal S_v(R)$.
\end{definition}
\begin{proof}
For each $r\in R$, the equality $|r|_x=|r|_v$ defines a closed subset
by continuity of evaluation and the Hausdorff property of $\R_{\ge0}$.
Their intersection is the relative spectrum, proving closedness.
An $R$-algebra homomorphism preserves the images of scalars, so restriction
preserves each displayed equality. The restricted map is continuous because
the unrestricted map is continuous and both relative spectra have the subspace topology.
The identity map $R\to R$ preserves every evaluation of $v$, yielding the final point.
\end{proof}

\begin{theorem}[Compatibility with trivial-field analytification {\cite[\S2.1]{FHX}}]
\label{spectrum:trivial-field-bridge}
\lean{EnvelopeProperty.Berk.toSpectrum,
EnvelopeProperty.Berk.relativeSpectrumHomeomorph,
EnvelopeProperty.Berk.relativeSpectrumHomeomorph_apply_val,
EnvelopeProperty.Berk.relativeSpectrumHomeomorph_comap}
\leanok
\uses{def:berk,spectrum:relative}
For a trivially valued field $k$ and a commutative $k$-algebra $A$, the affine
analytification $(\Spec A)^\an$ is canonically homeomorphic to
$\mathcal S_{\mathrm{triv}}(A)$. This homeomorphism preserves every seminorm evaluation
and commutes with restriction along $k$-algebra homomorphisms.
\end{theorem}
\begin{proof}
Both spaces consist of multiplicative ultrametric seminorms on $A$ whose value
on a nonzero scalar is one. The equality on the zero scalar is automatic because
seminorms and ring homomorphisms preserve zero. Thus the identity on seminorms
is a bijection between the two descriptions. Both topologies are the pointwise
evaluation topology, so the bijection and its inverse are continuous.
Restriction on both descriptions is composition of the same seminorm with the
same algebra homomorphism, proving naturality.
\end{proof}
